\subsection{Ablation Study}
Table~\ref{tab:ablation} reports the restart-related variants for which complete results are currently available. The rows of ARPSO-Fixed, ARPSO-Global, and ARPSO-Local are reserved for the six-variant ablation study and are left blank in this draft to avoid reporting unverified values.

\begin{table}[!t]
\caption{Ablation Results of Restart-related Variants}
\centering
\begin{tabular}{lccc}
\toprule
Algorithm & Avg. rank & Avg. restarts & Avg. runtime (s) \\
\midrule
PSO-RS & 6.07 & 8.94 & 76.21 \\
ARPSO-Fixed & -- & -- & -- \\
ARPSO-Global & -- & -- & -- \\
ARPSO-Local & -- & -- & -- \\
ARPSO-SRR & 4.31 & 67.72 & 71.82 \\
ARPSO-EIS & 4.45 & 66.55 & 72.07 \\
\bottomrule
\end{tabular}
\label{tab:ablation}
\end{table}

The completed restart-related results show that ARPSO-SRR obtains a better average rank than PSO-RS and is slightly better than ARPSO-EIS. The result indicates that the adaptive hybrid restart framework is more effective than a simple random restart strategy. The comparison with ARPSO-EIS also indicates that the more complex inefficient-particle identification rule does not provide stable additional improvement in the completed experiment.

Fig.~\ref{fig:ablation} is regenerated from the completed average-rank results. Only the completed restart-related variants are plotted to avoid showing uncompleted ablation variants as numerical results.

\begin{figure}[!t]
    \centering
    \safeincludegraphics[width=\linewidth]{ablation_bar_tikz_final.pdf}
    \caption{Average-rank comparison of completed restart-related variants.}
    \label{fig:ablation}
\end{figure}

